% =====================================================
% part9_applications.tex ― 第9部:応用展望
% =====================================================
\part{応用展望}

% =====================================================
\chapter{マルコフ連鎖とペロン・フロベニウスの定理}\label{chap:markov}
% =====================================================

\section{確率の行列}

「明日の状態は今日の状態だけで(確率的に)決まる」という型の現象
――天気、ブランド選択、ウェブの閲覧、塩基配列――を
\textbf{マルコフ連鎖}という。状態が $n$ 個あるとき、
状態 $j$ から状態 $i$ へ移る確率 $p_{ij}$ を並べた行列 $P$
(各\textbf{列}の和が $1$。これを\textbf{確率行列}という)を使うと、
時刻 $k$ の確率分布ベクトル $\vect{x}_k$ の時間発展は
\[
\vect{x}_{k+1} = P\,\vect{x}_k, \qquad \vect{x}_k = P^k \vect{x}_0
\]
と、\textbf{行列のべき乗}の問題になる。第IV部の主戦場である。

\begin{rei}
晴れ($1$)と雨($2$)の二状態。晴れの翌日は確率 $0.8$ で晴れ、
雨の翌日は確率 $0.4$ で晴れに回復するとする:
\[
P = \mat{0.8 & 0.4 \\ 0.2 & 0.6}.
\]
固有多項式は $\lambda^2 - 1.4\lambda + 0.4 = (\lambda - 1)(\lambda - 0.4)$。
固有値は $1$ と $0.4$ である。
\end{rei}

固有値に $1$ が現れたのは偶然ではない。各列の和が $1$ であることから、
行ベクトル $\mat{1 & 1 & \cdots & 1}$ が ${}^tP$ の固有値 $1$ の
固有ベクトルになり($A$ と ${}^tA$ の固有値は一致するので)、
\textbf{確率行列は必ず固有値 $1$ をもつ}。さらに、確率の総和が保存される
仕組みから、すべての固有値は $|\lambda| \le 1$ を満たす。

固有値 $1$ の固有ベクトルを、成分の和が $1$ になるように正規化したもの
$\vect{\pi}$($P\vect{\pi} = \vect{\pi}$)を\textbf{定常分布}という。
上の例では $\vect{\pi} = \mat{2/3 \\ 1/3}$ であり(確かめよ)、
任意の初期分布から出発しても
\[
\vect{x}_k = \vect{\pi} + (\text{固有値 } 0.4 \text{ の項}) \cdot 0.4^k
\longrightarrow \vect{\pi}
\]
と、定常分布に幾何級数的な速さで収束する。
\textbf{長期の運命は第二固有値が決め、第一固有値($=1$)の固有ベクトルが
行き先になる}――フィボナッチ数列(第IV部)・微分方程式(第VI部)で見た
支配原理の、確率版である。

\begin{teiri}[ペロン・フロベニウスの定理(概要)]
成分がすべて正の正方行列は、絶対値最大の固有値として
\textbf{正の実数の固有値をただ一つ}もち、それに属する固有ベクトルとして
\textbf{成分がすべて正のもの}が(定数倍を除き一意に)取れる。
\end{teiri}

この定理が、定常分布の存在・一意性と収束を保証する
(確率・人口・経済など「負の値をとらない量」の行列モデル全般の基礎定理である)。

\begin{itembox}[l]{深掘り:Google を生んだ固有ベクトル ― ページランク}
「重要なページとは、重要なページからリンクされているページである」
という再帰的な定義を、初期の Google は固有値問題として実装した。
ウェブページ $j$ から $i$ へのリンクを確率遷移とみなした巨大な確率行列
$P$(数十億次元!)を作り、その\textbf{定常分布 $\vect{\pi}$ の第 $i$ 成分}を
ページ $i$ の重要度(ページランク)とするのである。
ランダムにリンクをたどり続けるネットサーファーが、
長期的に滞在する確率の高いページほど重要、という解釈になる。
計算には次章で学ぶべき乗法(まさに $P^k\vect{x}_0$ の反復)が使われ、
ペロン・フロベニウスの定理が答えの存在と一意性を保証する。
\textbf{検索エンジンの心臓部は、巨大な行列の固有ベクトル計算}なのである。
\end{itembox}

\begin{mondai}\label{q:markov}
確率行列 $P = \mat{0.9 & 0.3 \\ 0.1 & 0.7}$ について、
固有値をすべて求め、定常分布 $\vect{\pi}$ を求めよ。
\end{mondai}

% =====================================================
\chapter{主成分分析 ― データの中の固有ベクトル}\label{chap:pca}
% =====================================================

\section{分散を最大にする方向}

$n$ 人について $2$ 科目の得点を測ったデータを、平均を引いて
(中心化して)偏差ベクトルの組にしたとする。
\textbf{主成分分析(PCA)}の問いはこうである:
「データが最も大きくばらつく方向(単位ベクトル $\vect{u}$)はどれか」。

方向 $\vect{u}$ にデータを正射影したときの分散は、
分散共分散行列 $S$(対称行列。第V部の深掘り参照)を使って
\[
(\vect{u} \text{ 方向の分散}) = {}^t\vect{u}\,S\,\vect{u}
\]
という\textbf{二次形式}で書ける。これを $\|\vect{u}\| = 1$ のもとで
最大化したい。スペクトル分解が一撃で答える。

\begin{teiri}
対称行列 $S$ の固有値を $\lambda_1 \ge \cdots \ge \lambda_n$、
対応する正規直交固有ベクトルを $\vect{u}_1, \dots, \vect{u}_n$ とすると、
$\|\vect{u}\| = 1$ の範囲で ${}^t\vect{u}S\vect{u}$ の最大値は
$\lambda_1$ であり、$\vect{u} = \vect{u}_1$ で達成される。
\end{teiri}

証明は鮮やかである。$\vect{u}$ を固有ベクトルの基底で
$\vect{u} = c_1\vect{u}_1 + \cdots + c_n\vect{u}_n$($\sum c_i^2 = 1$)と
展開すると
\[
{}^t\vect{u}S\vect{u} = \lambda_1 c_1^2 + \cdots + \lambda_n c_n^2
\le \lambda_1(c_1^2 + \cdots + c_n^2) = \lambda_1.
\]
等号は $c_1 = \pm 1$、すなわち $\vect{u} = \pm\vect{u}_1$ のとき。
\textbf{最大固有値の固有ベクトル $\vect{u}_1$ が第 $1$ 主成分}であり、
それと直交する範囲で分散最大の方向が $\vect{u}_2$(第 $2$ 主成分)、
以下同様である。全分散に占める $\lambda_i$ の割合
$\lambda_i / (\lambda_1 + \cdots + \lambda_n)$ を\textbf{寄与率}といい、
「上位 $k$ 主成分でデータの何割を説明できるか」の指標になる。

中心化したデータ行列に第VII部の SVD を施すと、
$V$ の列がそのまま主成分、特異値の二乗が(定数倍を除き)$\lambda_i$ になる。
\textbf{PCA とはデータ行列の SVD であり、上位 $k$ 主成分への要約とは
エッカート・ヤングの最良低ランク近似}である――第V部・第VII部の理論が、
ここで一つの実用手法に合流する。

\begin{itembox}[l]{深掘り:固有顔 ― 高次元データの圧縮の実際}
顔写真($100 \times 100$ 画素なら $1$ 枚 $= \R^{10000}$ のベクトル)を
大量に集めて PCA を行うと、上位の主成分はそれ自身が
「顔らしい画像」(\textbf{固有顔}と呼ばれる)になり、
個々の顔は $50$〜$100$ 個程度の主成分の線形結合で十分に近似できる
ことが知られている。$1$ 万次元のデータの実効的な次元が
$2$ 桁しかない――この「次元の劇的な圧縮可能性」は顔に限らず、
音声・文章・購買履歴など自然界のデータに広く見られ、
機械学習が機能する根本的な理由の一つだと考えられている。
現代の深層学習は PCA の線形写像を非線形写像に置き換えた
「表現学習」と見ることができ、その出発点に PCA がある。
なお、主成分の解釈(第 $1$ 主成分は「総合力」、第 $2$ は「文理の差」など)は
分析者の仕事であり、数学は方向を教えるだけである点には注意したい。
\end{itembox}

\begin{mondai}\label{q:pca}
中心化済みの $2$ 次元データ
$\mat{-1 \\ -1}, \mat{0 \\ 0}, \mat{1 \\ 1}$ について、
行列 $S = \sum_i \vect{d}_i\,{}^t\vect{d}_i$ を計算し、
第 $1$ 主成分の方向と寄与率を求めよ。
\end{mondai}

% =====================================================
\chapter{数値線形代数 ― 計算機の上の線形代数}\label{chap:numerical}
% =====================================================

\section{有限桁の世界}

コンピュータは実数をおよそ $16$ 桁の有限精度でしか持てない。
このため「数学的に正しい方法」が「数値的に良い方法」とは限らない。
本書でも既にいくつか出会った:行列式の定義($n!$ 項)は使わず掃き出し法を
使う(第II部)、固有値は固有方程式を解かず QR 法で求める(第VII部)、
正規方程式より QR 分解が安全(第VII部)。本章はこの世界の見取り図を与える。

\section{LU 分解 ― 掃き出し法の保存形式}

掃き出し法(前進消去)の操作を記録すると、$A$ は
\[
A = LU \qquad (L \text{ は下三角、} U \text{ は上三角})
\]
に分解される。一度 $LU$ を作っておけば、$A\vect{x} = \vect{b}$ は
$L\vect{y} = \vect{b}$(上から代入)と $U\vect{x} = \vect{y}$(下から代入)の
二度の代入で解ける。\textbf{右辺 $\vect{b}$ を取り替えて何度も解く}場面
(シミュレーションで頻出)では、分解を使い回せるこの形式が標準である。
逆行列をあらわに計算することは、実務ではほとんどない
($A^{-1}\vect{b}$ と書いてあったら「$LU$ 分解で解け」と読むのが慣わしである)。

\section{条件数 ― 問題そのものの敏感さ}

\begin{teigi}[条件数]
正則行列 $A$ の最大特異値と最小特異値の比
\[
\kappa(A) = \frac{\sigma_1}{\sigma_n}
\]
を $A$ の\textbf{条件数}という。
\end{teigi}

$A\vect{x} = \vect{b}$ において、入力 $\vect{b}$ の相対誤差は、
解 $\vect{x}$ では最大 $\kappa(A)$ 倍に拡大されうる
($A$ は最大 $\sigma_1$ 倍に伸ばし、$A^{-1}$ は最大 $1/\sigma_n$ 倍に
伸ばすため)。$\kappa$ が $10^{12}$ もあれば、$16$ 桁の計算でも
答えは $4$ 桁しか信用できない。重要なのは、これが\textbf{解法ではなく
問題自体の性質}だということである。条件数の悪い問題は、
どんな名アルゴリズムでも救えない――SVD が「問題の健康診断」の
道具にもなっている。

\section{べき乗法 ― 最大固有値だけ欲しいとき}

ページランクのような巨大行列では、全固有値は要らない。
最大のものだけ欲しい。そこで、適当な初期ベクトルに
$A$ を\textbf{ひたすら掛け続ける}:
\[
\vect{x}_{k+1} = \frac{A\vect{x}_k}{\|A\vect{x}_k\|}.
\]
固有ベクトル展開 $\vect{x}_0 = c_1\vect{v}_1 + \cdots + c_n\vect{v}_n$ に
$A^k$ を施すと $A^k\vect{x}_0 = c_1\lambda_1^k\vect{v}_1 + \cdots$ となり、
$|\lambda_1| > |\lambda_2|$ なら第 $1$ 項以外は相対的に
$(\lambda_2/\lambda_1)^k$ の速さで消えていく。
\textbf{「最大固有値が長期を支配する」という本書の支配原理を、
そのままアルゴリズムにした}のがべき乗法である。

\begin{itembox}[l]{深掘り:大規模疎行列の時代}
ウェブのリンク行列や SNS の友人関係行列は、次元こそ数十億だが、
成分のほとんどは $0$ である(\textbf{疎行列})。疎行列では
「行列を保存する」代わりに「零でない成分だけ保存する」ことで記憶量を、
「$A\vect{x}$ を高速に計算できる」ことを活かした反復法
(べき乗法はその最简単な例)で計算量を、それぞれ劇的に削減できる。
現代の数値線形代数は「$A$ の成分を見る」ことすら諦め、
「$\vect{x} \mapsto A\vect{x}$ という操作だけが与えられている」
という前提で設計される。深層学習の訓練も、本質的には
巨大な行列・テンソル演算の塊であり、GPU とはこの種の計算に
特化した装置である。\textbf{線形代数の計算効率が、そのまま
現代の計算機文明の速度を決めている}と言って誇張ではない。
\end{itembox}

\begin{mondai}\label{q:power2}
$A = \mat{2 & 1 \\ 1 & 2}$、$\vect{x}_0 = \mat{1 \\ 0}$ から出発して
$\vect{x}_1 = A\vect{x}_0$、$\vect{x}_2 = A\vect{x}_1$ を計算せよ
(正規化は省略してよい)。$\vect{x}_2$ の方向が、最大固有値 $3$ の
固有ベクトル $\mat{1 \\ 1}$ の方向に近づいていることを確認せよ。
\end{mondai}

% =====================================================
\chapter{線形代数の行き先 ― 経済・量子・AI}\label{chap:final}
% =====================================================

\section{産業連関分析 ― 経済全体を一つの行列に}

経済学者レオンチェフは、経済全体を一つの行列に書いた(ノーベル経済学賞)。
産業 $j$ が $1$ 単位生産するのに必要な産業 $i$ の財の量を $a_{ij}$ とした
\textbf{投入係数行列} $A$ を使うと、総生産 $\vect{x}$ のうち
$A\vect{x}$ は生産活動自身に食われ、残りが最終需要 $\vect{d}$ に回る:
\[
\vect{x} - A\vect{x} = \vect{d}
\quad\Longrightarrow\quad
\vect{x} = (E - A)^{-1}\vect{d}.
\]
「需要 $\vect{d}$ を満たすには、各産業がいくら生産すべきか」が
逆行列一発で求まる。さらに $(E - A)^{-1}$ は、等比級数
$\dfrac{1}{1 - a} = 1 + a + a^2 + \cdots$ の行列版
\[
(E - A)^{-1} = E + A + A^2 + A^3 + \cdots
\]
で書ける($A$ の固有値の絶対値がすべて $1$ 未満のとき収束。
$e^A$ と同じく「級数に行列を代入する」発想である)。各項の経済学的な
意味は鮮明で、$E$ は需要そのもの、$A\vect{d}$ はそれを作るための直接の
中間需要、$A^2\vect{d}$ はそのまた中間需要…と、
\textbf{波及効果の無限連鎖}を表す。公共事業の経済波及効果の試算などで、
この行列(レオンチェフ逆行列)は現在も実務で使われている。
収束の保証にはペロン・フロベニウスの定理(投入係数は非負!)が働いている。

\begin{mondai}\label{q:leontief}
投入係数行列 $A = \mat{0.2 & 0.4 \\ 0.3 & 0.1}$、最終需要
$\vect{d} = \mat{6 \\ 3}$ のとき、必要な総生産 $\vect{x} = (E - A)^{-1}\vect{d}$
を求めよ。また、得られた $\vect{x}$ が $\vect{x} - A\vect{x} = \vect{d}$ を
満たすことを検算せよ。
\end{mondai}

\section{量子力学 ― 線形代数が世界の文法である世界}

量子力学の数学的骨格は、本書の語彙でほぼ完全に記述できる。
状態は(複素)ベクトル空間の単位ベクトル、物理量は対称行列の複素版である
\textbf{エルミート行列}、観測で得られる値はその\textbf{固有値}、
観測直後の状態は対応する固有空間への\textbf{正射影}(スペクトル分解!)、
時間発展は $e^{-itH}$ という\textbf{行列の指数関数}(直交行列の複素版
=ユニタリ行列になる)、複数の粒子の合成系は\textbf{テンソル積}、
もつれは\textbf{分解不能テンソル}。
――第V〜VIII部で学んだ概念が、一つ残らず物理の最深部の語彙として
再登場する。「自然はなぜ線形代数で書かれているのか」は誰にも分からないが、
書かれていること自体は $20$ 世紀物理学の最大の発見の一つである。

\section{AI ― 行列計算の文明}

現代の深層学習・大規模言語モデルの内部は、誇張なく行列演算の積み重ねである。
ニューラルネットの一層は「行列を掛けて、非線形関数を通す」操作であり、
学習とはその行列(重み)を勾配で更新し続けることである。
単語や画像は高次元ベクトル(\textbf{埋め込み})として表現され、
意味の近さが内積で測られ(第V部)、注意機構(Attention)は
内積による類似度行列の計算を核とし、
モデルの圧縮には低ランク近似(第VII部)が使われる。
本書を理解した読者は、AI の解説論文に現れる数式の大部分を、
既に「読める」はずである。

\section*{おわりに ― 二本の柱を振り返る}

本書は「線形代数とは、連立一次方程式と線形変換を統一的に扱う数学である」
という一本の柱から出発した。長い旅を終えたいま、
全体は次のように要約できる。

\begin{itembox}[l]{本書の全景}
\begin{itemize}
\item \textbf{解く}(第I・II部):掃き出し法・行列式・階数が、
連立方程式の三択(一意・無数・なし)を完全に裁定した。
\item \textbf{抽象化する}(第III部):公理化により、多項式も関数も数列も
同じ舞台に乗った。
\item \textbf{分解する}(第IV〜VII部):固有値・対角化・ジョルダン標準形・
スペクトル分解・SVD。「よい基底を選べば、どんな変換も
伸縮の重ね合わせに見える」。べき乗・指数関数・微分方程式・最小二乗が
その系として落ちてきた。
\item \textbf{構成する}(第VIII部):双対・商・テンソル・外積。
「線形」という概念そのものの文法を手に入れ、行列式の正体を見届けた。
\item \textbf{使う}(第IX部):確率・データ・計算・経済・量子・AI。
支配原理(最大固有値が長期を決める)と分解の思想が、
現実世界を動かしていた。
\end{itemize}
\end{itembox}

第I部で $2$ 本だった正則性の同値条件は、いまや
「$\det A \neq 0$、$\mathrm{rank}\,A = n$、列が線形独立、
$\mathrm{Ker} = \{\vect{0}\}$、$0$ が固有値でない、
特異値がすべて正…」と長大なリストに成長した。
\textbf{一つの概念を多くの言葉で言い換えられること、
一つの定理が多くの分野で同じ顔をして現れること}
――線形代数の見通しのよさの正体は、この往復運動にある。

この先に進む読者へ。解析方向なら\textbf{関数解析}
(無限次元のスペクトル理論。量子力学の本来の住処)、
代数方向なら\textbf{表現論}(群の対称性を行列で写し取る理論)と
\textbf{加群の理論}(ジョルダン標準形の真の証明が住む場所)、
計算方向なら\textbf{数値線形代数}と\textbf{最適化}、
データ方向なら\textbf{統計的機械学習}が、それぞれ本書の続きである。
どの道を選んでも、基底を取り替え、分解し、不変量を見つけるという
本書の身のこなしは、そのまま通用する。

行列の積の奇妙な定義に戸惑った第I部の自分を、思い出してほしい。
定義の必然性が見えたとき、数学は暗記から物語に変わる
――その経験を、これからも積み重ねていってほしい。
